\chapter*{Abstract}
\phantomsection
\addcontentsline{toc}{chapter}{Abstract}
\markboth{Abstract}{Abstract}

Harmony Hub is a mobile application designed to transform a user's emotional and
contextual state into a more useful musical experience for everyday wellbeing.
The system combines a guided check-in, contextual recommendation logic, a
structured music catalogue, Spotify-based playlist materialization and a
feedback-based learning cycle. From a technical point of view, the solution is
built around a Flutter mobile client, a FastAPI backend, document persistence
in Firestore and a supervised session-level model that is only enabled when
there is sufficient objective evidence to justify its use.

The recommendation process does not rely exclusively on a machine learning
model. Instead, it follows a hybrid approach: a heuristic layer guarantees
operability from the first session, while a logistic regression model can later
reorder candidate session modes when enough historical data, class balance and
quality metrics are available. The musical ranking itself is primarily carried
out over the project's own \texttt{msd\_tracks} catalogue, whereas Spotify is
used mainly for authorization, profile context and final playlist
materialization. This design allows the application to remain useful even under
external constraints such as limited Spotify endpoints or insufficient training
data. The result is not merely a musical suggestion, but a real playlist
generated in the user's Spotify account and linked to a traceable session
history.

The project therefore explores how mobile software, external APIs, contextual
data and progressive learning can be combined to support everyday moments of
emotional regulation without positioning the app as a clinical tool. In its
current state, the user-level progressive profile is already operational and
visible in the experience, with longitudinal evidence across the five
emotional states considered in the validation process. The supervised
session-level model is also active under explicit quality gates and can
contribute probabilistic reordering of candidate session modes while still
preserving abstention when evidence is weak. The validation additionally shows
that environmental sensing can be incorporated into both the recommendation
context and the playlist generation pipeline. The outcome is a functional
prototype that demonstrates the viability of a wellness-oriented,
session-based and explainable music recommendation approach with real Spotify
playlist creation, contextual traceability and progressive learning.
